% !TeX spellcheck = ru_RU
% !TEX root = vkr.tex

\section*{Введение}
\thispagestyle{withCompileDate}

Морфометрия~--- количественное описание формы биологических объектов~--- является одним из ключевых инструментов таксономии и экологии~\cite{Bookstein1991,DrydenMardia1998}.
Несмотря на развитый аппарат статистического анализа, получение исходных данных по-прежнему остаётся ручной или полуавтоматической операцией: исследователь вручную обводит контур каждого объекта в растровом редакторе и вручную снимает измерения.
При сериях из десятков и сотен образцов это занимает значительное время и вносит субъективную погрешность, особенно когда измерения ведут несколько человек~\cite{Bookstein1998}.

Задача осложняется при работе с морфологически сложными объектами.
Листья мхов рода \textit{Sphagnum}~--- экологически значимых эдификаторов торфяных болот и тундровых сообществ~\cite{ClymoHayward1982,Rydin2006}~--- состоят из одного слоя клеток, могут быть прозрачными и иметь неоднородные, рваные края.
Существующие программы автоматической «обрисовки» контуров (например, tpsDIG~\cite{Rohlf2021}) требуют хорошо контрастированных объектов без инородных включений и не справляются с такими объектами.

С технической точки зрения проблема состоит в отсутствии инструмента, объединяющего в едином рабочем цикле нейросетевую сегментацию, интерактивную коррекцию результата и автоматическое извлечение морфометрических параметров с экспортом в табличный формат.
Появление foundation-моделей сегментации~--- в частности, Segment Anything Model (SAM)~\cite{Kirillov2023} и её облегчённой версии MobileSAM~\cite{Zhang2023}~--- открывает возможность строить такой инструмент без необходимости собирать размеченные обучающие данные.
Настоящая работа направлена на решение описанной проблемы.

Работа имеет следующую структуру.
В разделе~\ref{sec:task} сформулированы цель, задачи и требования к разрабатываемой системе.
В разделе~\ref{sec:relatedworks} проведён обзор существующих инструментов и обоснован выбор модели сегментации.
В разделе~\ref{sec:solution} описаны архитектура и реализация приложения.
В разделе~\ref{sec:experiment} приведены результаты сравнения с ручной разметкой и апробации.
Заключение содержит выводы и направления дальнейшего развития системы.
